\documentclass[11pt,a4paper]{article}
\usepackage[T1]{fontenc}
\usepackage[utf8]{inputenc}
\usepackage[margin=17mm]{geometry}
\usepackage{amsmath,amssymb}
\usepackage{mathptmx}
\usepackage{microtype}
\usepackage{titlesec}
\usepackage[hidelinks]{hyperref}
\linespread{0.96}
\hypersetup{pdftitle={Homotopy methods for optimization},
pdfauthor={Aubin Charley, Alexandre Corrard, Idriss El khamlichi, Maxime Nicaise}}

\titleformat{\section}{\normalsize\bfseries}{\thesection.}{0.5em}{}
\titlespacing*{\section}{0pt}{8pt plus 1pt minus 1pt}{4pt}
\setlength{\parindent}{1em}
\setlength{\parskip}{2pt plus 0.5pt minus 0.5pt}
\setlength{\abovedisplayskip}{5pt plus 1pt minus 1pt}
\setlength{\belowdisplayskip}{5pt plus 1pt minus 1pt}
\setlength{\abovedisplayshortskip}{3pt}
\setlength{\belowdisplayshortskip}{3pt}
\setlength{\footskip}{8mm}
\raggedbottom

\begin{document}
\begin{center}
    {\Large\bfseries Homotopy methods for optimization}\par
    \vspace{4pt}
    Aubin Charley \quad Alexandre Corrard \quad
    Idriss El khamlichi \quad Maxime Nicaise\par
    \vspace{2pt}
    September 2026
\end{center}
\vspace{-4pt}

An idea to minimize a difficult target objective $L(\theta)$ is to introduce
an approximation family $(L_s)_{s\in[0,1]}$, with $L_1=L$, such that the
initial problems are easier to solve. The family progressively evolves to
the target problem. In practice, we consider $0=s_0<\cdots<s_K=1$ and use
an approximate minimizer of $L_{s_k}$ as a warm start for $L_{s_{k+1}}$:
$\theta_{k+1}^{(0)}=\theta_k^\star$. The best-case scenario is that each
solution lands in a basin where a local optimizer finds the next minimizer.
However, the central difficulty is to explain why the intermediate problems
should be easier, and why solving them should help with the final one.
The approaches below give different answers to these two questions.

\section{Making the objective easier to optimize}

The most intuitive approach is to progressively deform an easy objective
into the target loss, for instance through
\[
    L_s(\theta)=(1-s)L_{\mathrm{easy}}(\theta)+sL(\theta).
\]
In \cite{LENDL1998359}, the authors apply the deformation principle to
multilayer perceptrons, starting from a modified regression problem with
a known solution. The main difficulty is the choice of $L_{\mathrm{easy}}$:
there is no general reason for the intermediate objectives to remain easier
to optimize, or for their minimizers and attraction basins to vary smoothly
along the path. A related construction uses regularization,
$L_s(\theta)=L(\theta)+\lambda(s)R(\theta)$, with $\lambda(s)\downarrow0$.
In \cite{amakor2025continuation}, the authors study regularization paths for
deep networks by treating the loss and $\ell^1$ sparsity as competing
objectives. Their aim is to explore this trade-off; strong regularization
does not by itself guarantee an easy starting problem.

Parameter-space smoothing gives a more direct geometric motivation:
the loss is averaged around each parameter value, and the amount of
smoothing is progressively reduced. Writing
\[
    L_\sigma(\theta)=\mathbb{E}_{\xi}
    \bigl[L(\theta+\sigma\xi)\bigr],
    \qquad \sigma\downarrow0,
\]
In \cite{hazan2016graduated}, the authors consider $\xi$ uniform on the unit
ball. Under assumptions ensuring local strong convexity of the smoothed
objectives around their minimizers and sufficient proximity of successive
minimizers, they obtain convergence guarantees for graduated optimization.
This makes explicit what the interpolation alone leaves unresolved:
smoothing must provide both tractable subproblems and a path that can be
followed. The approach looks like a good idea from a purely mathematical
perspective, because smoothing acts directly in the optimization space.
However, these geometric assumptions are difficult to verify for
neural-network losses. In \cite{sato2025explicit}, the authors investigate
this transfer to modern networks and report limitations of explicit
smoothing on high-dimensional objectives. They also study an implicit
alternative, interpreting SGD stochasticity as a source of smoothing,
and extend it to SGD with momentum. Thus even a well-motivated family
raises the practical question of how much its intermediate problems cost
to optimize.

\section{From geometric simplicity to simpler examples and models}

The previous constructions act directly on the objective. Curriculum
learning, introduced in \cite{bengio2009curriculum}, instead aims at building training
losses that become progressively more difficult from a semantic point of
view. It gives more importance to intuitively ``easy'' examples at first,
before recovering the target training distribution:
\[
    L_s(\theta)=\sum_{i=1}^{n}w_i(s)\,
    \ell\bigl(f_\theta(x_i),y_i\bigr),
    \qquad w_i(s)\geq0,\quad \sum_i w_i(s)=1,\quad w_i(1)=\tfrac1n.
\]
In \cite{bengio2009curriculum}, the authors show, for instance, that
presenting gradually more complex geometrical shapes can help models
learn faster. The main limitation is that the notion of an ``easy''
example is highly task-dependent and generally defined from prior
knowledge. Moreover, semantic simplicity does not imply a simpler
optimization landscape. This is the main gap between the curriculum
intuition and the geometric conditions used in \cite{hazan2016graduated}.

A similar intuition can be applied inside the network. If
$f_{\theta,s}$ denotes a progressively modified model, the family becomes
$L_s(\theta)=n^{-1}\sum_i\ell(f_{\theta,s}(x_i),y_i)$, with
$f_{\theta,1}=f_\theta$. In \cite{gulcehre2017mollifying}, the idea is to
make the network almost linear at the beginning of training through
stochastic shortcuts and noisy activations, then progressively remove
these modifications. In \cite{yang2025homotopy}, the authors use an
activation homotopy connecting the identity to ReLU,
$\phi_s(z)=(1-s)z+s\operatorname{ReLU}(z)$, and obtain accelerated
convergence in an infinite-width two-layer setting. In
\cite{roulet2021smoothing}, the authors instead replace non-smooth network
modules by smooth approximations. Their fixed-smoothing analysis provides
a possible family for continuation, but is not itself a homotopy-training
algorithm.

These constructions simplify the model in different senses;
in particular, a linear input-output map still need not give a convex
loss in the parameters of a multilayer network.

Frequency-based curriculum connects this model modification to the idea
of gradually increasing the difficulty of the information presented.
In \cite{sinha2020curriculum}, the authors initially remove high-frequency
information from CNN feature maps and progressively reintroduce it:
\[
    h_{\ell,\sigma}=G_\sigma*h_\ell,
    \qquad \sigma\downarrow0,
\]
where $G_\sigma$ is a Gaussian low-pass filter with standard deviation
$\sigma$. The intuition is that the model can first focus on large-scale
structures before dealing with finer details and textures. The spectral
bias studied in \cite{rahaman2019spectral}, whereby
networks tend to learn low-frequency components first, gives a
conceptually related motivation. However, the two notions of frequency
are not the same: \cite{rahaman2019spectral} studies the learned function
in input space, while \cite{sinha2020curriculum} directly filters spatial feature maps. This supports a useful analogy,
but does not establish that feature smoothing satisfies the geometric
conditions of graduated optimization.

\section{Once a family is chosen: tracking or learning the path}

The choice of a family does not determine how efficiently its solutions
can be reached. In \cite{mai2026neural}, the authors address this second
question with a neural predictor-corrector: a reinforcement-learning
policy adapts the continuation step size and the termination of corrector
iterations. The policy can learn to take larger steps when the path is
easy to track and smaller steps in more difficult regions, while adapting
the computation spent on each correction. This reduces manual schedule
tuning, but requires training on a suitable distribution of problems.

In \cite{Lin2023ContinuationPL}, the authors go further by learning the
solution path itself rather than explicitly tracking it through a sequence
of optimization problems. A network $\theta_\phi(s)$ represents the
solutions across the continuation interval, schematically through
\[
    \min_\phi\ \mathbb{E}_{s\sim\mathcal{U}(0,1)}
    \bigl[L_s\bigl(\theta_\phi(s)\bigr)\bigr].
\]
This replaces the cost of sequential path tracking by the cost of learning
a global representation, which can itself be difficult when minimizers
undergo sharp transitions or form competing branches. Both approaches
still require a suitable approximation family to be defined in advance:
learning how to follow or represent a path does not settle how to
construct a useful one.

For our exploratory project, this suggests separating the benefit of the
chosen simplification from the benefit of introducing difficulty
progressively. Parameter smoothing offers the clearest connection to
optimization geometry, while activation and feature smoothing offer
concrete constructions for neural networks. Comparing a decreasing
smoothing schedule with both fixed smoothing and direct training, under
a comparable computational budget, would help determine whether
continuation itself improves optimization. The literature provides
plausible starting points, but no general guarantee that the most
intuitive simplification will produce the most useful path.

\bibliographystyle{plain}
\bibliography{biblio}
\end{document}
